\begin{definition}
    $E_1, E_2$~--- ЛНП, $A \in \L(E_1, E_2)$. $A$ называется \textit{компактным}, если $\forall$ огр. $S \subset E_1$ $A[S]$~--- предкомпакт (в $E_2$). $K(E_1, E_2)$~--- множество компактных операторов
\end{definition}

\noindent\textbf{Напоминание.}
В банаховом пространстве предкомпактность эквивалентна вполне ограниченности

\begin{definition}[эквивалентное, <<с волной>>] % будем использовать
$E_1, E_2$~--- ЛНП, $A  \in \L(E_1, E_2)$ называется \textit{компактным}, если 
$A\left[ \overline{B_1}(0) \right]$~--- предкомпакт.

\textit{Аналогия с определением ограниченного оператора}
\end{definition}

\begin{proof}[эквивалентность]
Пусть у нас есть ограниченное множество $S$ и есть наш единичный шар $\overline{B}_1(0)$. Рассмотрим увеличинный в $R$ раз шар $R\overline{B}_1(0)$, который покроет $S$. Заметим, что если $A\left[\overline{B}_1(0) \right]$ --- предкомпакт, то в увеличенном виде он тоже будет вполне ограниченным множеством (у нас просто будет $R\varepsilon\text{-сеть вместо } \varepsilon\text{-сети}$).

Тогда $A[S] \subset A\left[R\overline{B}_1(0) \right]$. Так как $A\left[ R\overline{B}_1(0) \right]$ --- предкомпакт, то $A[S]$ тоже предкомпакт (как его подмножество).

Это следует из факта, что подмножество компактного множества --- предкомпакт.
\end{proof}

\begin{unnecessary}
\begin{definition}[Вполне непрерывные операторы (<<улучшающая сходимость>>)]
\[
x_n \weakconv x \Rightarrow \|Ax_n - Ax\| \to 0
\]
\end{definition}

\begin{exercise}
    $E$~--- банахово, $A \in K(E), E \ni x_n \weakconv x \in E \implies \|Ax_n - Ax\| \to 0$
    
    Подсказка:

    \[
    \begin{cases}
    y_n \weakconv y\\
    \{y_n\} \text{~--- предкомпакт}
    \end{cases}
    \Rightarrow \|y_n - y\| \to 0
    \]
\end{exercise}

\begin{lemmanote}
    Если $E$~--- рефлексивное банахово, то компактность $\sim$ вполне непрерывность.
\end{lemmanote}

\begin{example}
Ищем первообразную для функции $f(x)$:
$y' = f(x),\: y(0) = 0$. Это эквивалентно $y(x)=\int_{t=0}^x f(t)dt$

Резольвента~--- это <<решатель>> задачи. В данном случае <<решателем>> является оператор Вольтерра.

Вообще тут ситуация такая: у нас есть $C[0,1]$. На нем есть многообразие тех функций, которые в нуле равны нулю. На этом многообразии определен оператор дифференцирования, который дает нам все $C[0, 1]$, причём засчёт нормировки отображение взаимно однозначно, и есть обратный оператор. Он компактный, вообще хороший, всем советуем.
\end{example}
\end{unnecessary}

\begin{statement}[1]
Если $\dim E_1 < \infty$ или $\dim E_2 < \infty$,
$A \in \L(E_1, E_2)$, то $A$~--- компактный оператор
\end{statement}

\begin{proof}
\begin{itemize}
    \item $\dim E_1 < \infty$.
        $M \subset E_1$ ограничено, а значит вполне ограничено в конечномерном пространстве. $A$~--- непрерывный, а значит для
        $M \subset \overline{M}$~--- компакт, $A[M] \subset A[\overline{M}]$~--- компакт (образ компакта под действием непрерывного отображения~--- компакт). $\overline{A[M]} \subset A[\overline{M}]$.
        $\overline{A[M]}$~--- это замкнутое множество внутри компакта, то есть компакт, а значит $A[M]$ является предкомпактом. 
    \item $\dim E_2 < \infty$. $M \subset E_1$ ограничено, $A[M]$ ограничено в $E_2$, а значит вполне ограничено и является предкомпактом.
\end{itemize}
\end{proof}

\begin{unnecessary}
\begin{example}[В бесконечномерном пространстве всё сложнее]
$E = l_2(\C)$, $I$~--- тождественный оператор.

$I(\overline{B_1}(0)) = \overline{B_1}(0)$~--- не вполне ограничено.
\end{example}
\end{unnecessary}

\begin{statement}\label{6.1_st2}
    $K(E)$~--- идеал в $\L(E)$, то есть $\forall B \in \L(E) \: \forall A \in K(E) \implies AB, BA \in K(E)$. (Идеал в обычном алгебраическом смысле).
\end{statement}

\begin{proof}
Возьмём множество $M$, тогда $BM$ даст нам ограниченное множество в силу того, что $B$ ---ограниченный оператор. После применения $A$ получаем предкомпакт.

Если в другом порядке: $AM$ предкомпакт. Осталось показать, что после применения $B$ предкомпактность сохранится. Это действительно так, поскольку если взять $\varepsilon$-сеть в $AM$, и под действием оператора $B$ она превратится в $\|B\|\varepsilon$-сеть.
\end{proof}


\begin{statement}\label{6.1_st3}
Если $\dim E = \infty$, то $I \not \in K(E)$ ($I$~--- тождественный оператор).
\end{statement}

\begin{note}
Хотим узнать, конечномерно ли пространство $E$?
Теперь появился новый способ описания конечномерных пространств:
\[\dim E < \infty \Leftrightarrow \text{единичная сфера вполне ограничена} \Leftrightarrow I\text{~--- компактный оператор}\]
\end{note}

\begin{proof}
    Мы знаем, что при $\dim E < \infty \Leftrightarrow$ единичная сфера вполне ограниченное множество. А это уже равносильно тому, что $I\text{~--- компактный оператор}$.
\end{proof}

\begin{corollary}[Из утверждений~\ref{6.1_st2} и~\ref{6.1_st3}]
    $\dim E = \infty,\: A \in K(E) \implies 0 \in \sigma(A)$
\end{corollary}

\begin{proof}
    Предположим противное, то есть $\underbrace{(A - 0 I)^{-1}}_{=A^{-1}} \in \L(E)$. Но тогда в предположении, что 
$A$ -- компактный, можно получить $AA^{-1} = I$~--- компактный. Из того, что $\dim E = \infty$, следует, что $0$ лежит только в спектре.
\end{proof}

\begin{theorem}[12.1]\label{12.1}
    Пусть $E_1$~--- нормированное, $E_2$~--- банахово пространство, $\forall n \: A_n \in K(E_1, E_2)$ и $\|A_n - A\| \to 0$. Утверждается, что $A \in K(E_1, E_2)$.
\end{theorem}

%Не всегда компактные операторы представляются как предел операторов конечного ранга.

\begin{proof}
    Подсказка самим себе: аналогия с равномерной сходимостью непрерывных функций

\begin{itemize}
    \item[1.] Нужно доказать, что $A[\overline{B}]$~--- вполне ограниченное множество ($\overline{B}$ - единичный шар из определения <<с волной>>). Зафиксируем эпсилон из определения предела
    \[ \forall \varepsilon > 0 \: \exists N\!: \, \forall n \geqslant N \: \|A_n-A\| < \varepsilon\]
    % обводим себя в кружок чтобы мы никуда не убегали
    \item[2.] Фиксируем $n \geqslant N$
    $A_n[\overline{B}]$~--- вполне ограничено, следовательно
    $\exists \varepsilon$-сеть $\{y_1^n, \ldots y_{m_n}^n\}$ для $A_n[\overline{B}]$.
    \[
    \|A x - y_k^n\| \leqslant \|A x - A_n x\| + \|A_n x - y_k^n\| \leqslant 2 \varepsilon
    \]
\end{itemize}
\end{proof}

\begin{unnecessary}
    \begin{theorem}[Шаудера]
    $E$~--- банахово. $A$~--- компактный оператор тогда и только тогда, когда $A^*$~--- компактный оператор.
    \end{theorem}
    
    %нам вообще нужно доказательство теоремы Шаудера?
    % \begin{proof}
    % <<Я не знаю, как это доказать>> (c) С. П. Коновалов.
    % Докажите для случая $E$~--- банахово и линейных ограниченных операторов.
    % \end{proof}
\end{unnecessary}
